\documentclass[11pt,a4paper]{article}

\usepackage[margin=1.1in]{geometry}
\usepackage{amsmath,amssymb,amsthm,mathtools}
\usepackage[T1]{fontenc}
\usepackage{palatino}
\usepackage{microtype}
\microtypesetup{expansion=false}
\usepackage{tikz}
\usetikzlibrary{positioning,arrows.meta,calc,fit,backgrounds,patterns}
\usepackage{array}
\usepackage{booktabs}
\usepackage{longtable}
\usepackage{listings}
\usepackage{stmaryrd}
\usepackage[hidelinks]{hyperref}
\usepackage{enumitem}

\setlength{\emergencystretch}{2.5em}

\theoremstyle{plain}
\newtheorem{theorem}{Theorem}[section]
\newtheorem{proposition}[theorem]{Proposition}
\newtheorem{lemma}[theorem]{Lemma}
\newtheorem{corollary}[theorem]{Corollary}
\newtheorem{conjecture}[theorem]{Conjecture}
\newtheorem{principle}[theorem]{Principle}
\theoremstyle{definition}
\newtheorem{definition}[theorem]{Definition}
\newtheorem{observation}[theorem]{Observation}
\newtheorem{requirement}[theorem]{Requirement}
\newtheorem{example}[theorem]{Example}
\theoremstyle{remark}
\newtheorem{remark}[theorem]{Remark}

\newcommand{\rhoc}{\textsf{rho}}
\newcommand{\flang}{\textsf{f1r3lang}}
\newcommand{\Nil}{\mathbf{0}}
\newcommand{\sepc}{\mathbin{\big|}}
\newcommand{\out}{\mathsf{out}}
\newcommand{\inn}{\mathsf{in}}
\newcommand{\Rnn}{\mathbb{R}_{\geq 0}}
\newcommand{\Trop}{\mathbb{T}}
\newcommand{\Nms}{\mathcal{N}}
\newcommand{\Opn}{\mathcal{O}}
\newcommand{\dbr}[1]{\llbracket #1 \rrbracket}
\newcommand{\qt}[1]{\ensuremath{@#1}}
\newcommand{\Kmod}[1]{\langle K_{#1}\rangle}
\newcommand{\KKmod}[1]{\langle\!\langle K_{#1}\rangle\!\rangle}
\newcommand{\Bmod}[1]{[K_{#1}]}
\newcommand{\Emod}[2]{\langle E\rangle^{#2}_{#1}}
\newcommand{\gsep}[2]{\mathbin{|_{#1}^{#2}}}

\lstdefinestyle{rho}{
  basicstyle=\ttfamily\footnotesize,
  keywordstyle=\bfseries,
  commentstyle=\itshape,
  morekeywords={new,in,for,contract,match,if,else,true,false,Nil,where,and,or,not,
                exists,forall,nu,rules,rewrites,equations,types,language,name},
  morecomment=[l]{//},
  columns=fullflexible,
  breaklines=true,
  breakatwhitespace=false,
  showstringspaces=false,
  frame=none,
  xleftmargin=1em
}
\lstset{style=rho}

\title{\textbf{Sixteen Scientists on One Board}\\[0.4em]
\large A society of piece-agents as a chess player, in which coordination is\\
quorum sensing, strategy is a coalition formula, and a ply is a join}
\author{A F1R3FLY.io working note}
\date{September 2026}

\begin{document}
\maketitle

\begin{abstract}
This note specifies a chess player whose \emph{player} is not an evaluator but a
society: each piece of the moving colour is an agent in its own right, and the
sixteen agents coordinate to emit one move per turn. The board is likewise an
agent. The setting is cost-accounted \rhoc\ as implemented on the
\texttt{cost-accounted-rho} branch of \texttt{f1r3node-rust}, with the
generated spatial--behavioural logic of the mortal-scientist programme, and we
assume at development time an implementation of coalition logic augmented with
context-labelled modalities. We argue five structural points and then give a
spec. First, the board needs three name families per square, not one: an
occupancy cell, a move interface, and an attack-marker family, because
perception must not be able to act and because the attack map is the one
expensive object all sixteen agents share. Second, a chess ply is not a single
rewrite but an \emph{$N$-ary join}, and the join arity is exactly the move
taxonomy: two for a quiet move, two with occupant termination for a capture,
three for en passant, four for castling. This is what preserves the generated
modality $\Kmod{j}$ and the bijection between redex positions and legal moves.
Third, the coordination problem is the honeybee problem, and its solution is
forced by the same altitude argument: consensus over sixteen pieces is a global
universal with no presented witness, quorum at a candidate move is a local
$\oplus\otimes$ over one open, so the society must quorum-sense. Fourth --- and
this is the load-bearing observation --- in \rhoc\ a quorum \emph{is} a join
arity and one linear advocacy token per piece per round is the polity's
franchise, so the two-quorum theorem becomes an arithmetic fact about sends
rather than a simulation result: with sixteen pieces and threshold $k \geq 9$ a
split is structurally impossible. Fifth, legality is a genuine universal and
therefore belongs to the world, not to the society; pins and checks are
deposited structure and pin-awareness is perception. We then observe that
capture is predation in the existing sense, so material advantage is
metabolic advantage and the classical piece values become something the ecology
might derive rather than receive. The implementation spec that follows fixes
the namespace discipline, the turn protocol, the formula ladder, the cost
schedule, the requirements placed on the assumed coalition-logic checker, a six
phase delivery plan, and eight measurements, of which the sharpest asks whether
the noughts-and-crosses finding that non-modal rungs deliver ninety-two percent
of yield for thirty-four percent of cost is a fact about games or a fact about
noughts.
\end{abstract}

\tableofcontents

\section{Introduction}
\label{sec:intro}

A conventional chess program is a single agent. It holds a position, generates
moves, evaluates leaves, and returns an argmax. Everything about its
architecture follows from that: the evaluation function is global, the search is
centralised, and the pieces are data.

The proposal this note specifies inverts that. Each piece of the moving colour
is an agent with its own budget, its own hypothesis about the position, and its
own capacity to act. The board is an agent too. Sixteen agents want to move and
exactly one move commits, so the architecture is obliged to contain something
the conventional design never needs: a mechanism that turns sixteen concurrent
desires into one commitment, without a chairman.

That obligation is the reason to build the thing. It is not a distributed
implementation of a familiar algorithm. It is the smallest classical domain in
which the constructions of the mortal-scientist programme --- cost-accounted
computation, the generated spatial--behavioural logic, the assay in guard
position, predation as the sole route to budget, and coalition logic read as a
hypothesis language about populations --- all have non-trivial content at the
same time, and in which a century of external calibration is available to say
whether they are worth anything.

\subsection{What is being claimed, and what is not}

We are not claiming that a society of piece-agents will play strong chess. On
the contrary: \S\ref{sec:risks} argues it will play weak chess for a long time,
for reasons that are structural rather than incidental, and the note is written
so that this is not a surprise. What is claimed is that the society is a better
instrument than a conventional engine for a specific question the programme has
already asked and answered only in miniature, which is whether the yield of a
hypothesis is worth its price; and that chess is the smallest board on which the
answer could plausibly differ from the answer noughts-and-crosses gave.

\subsection{The spectrum, and where the interesting structure is}

The design space is usually presented as a spectrum. At one extreme \flang\
coordinates sixteen sessions with different language models, so a piece-agent is
a thin wrapper around a model session. At the other extreme each agent is
\flang\ logic throughout. Section~\ref{sec:policy} argues that the extremes are
not really two points on one dial. They are two answers to a single question ---
where a piece's policy lives, and whether it is priced --- and that the
programme already has a construction that lets one have both without incoherence,
namely the aperture. The recommendation is not a compromise between the extremes
but a relocation: the model goes outside the ply loop, where it is good at the
jobs it is good at, and inside only as a declared aperture in a stratified shard.

\subsection{Organisation}

Section~\ref{sec:imported} states what is imported from the existing programme
and what is assumed to exist at development time. Section~\ref{sec:why} argues
for chess as the next worked example. Sections \ref{sec:two}--\ref{sec:material}
are the analysis: two societies, the board, the ply as a join, the coordination
problem, coalition logic, and material as metabolism. Section~\ref{sec:policy}
handles the model axis. Section~\ref{sec:spec} is the implementation spec and is
the longest section. Section~\ref{sec:measure} lists the measurements.
Section~\ref{sec:risks} is the honest accounting of what will go wrong.
Section~\ref{sec:open} lists open problems. Appendix~\ref{app:rho} is the
\rhoc\ listing and Appendix~\ref{app:formulae} the formula catalogue.

\section{What is imported, and what is assumed}
\label{sec:imported}

\subsection{From the mortal-scientist programme}

The note stands on five constructions that already exist and are not re-derived
here.

\paragraph{The cut and the three grades of access.} A computation is
ontologically isolated; the environment of a computation is other computations;
and a learner reaches its environment by perception (reading deposited
structure, priced by formula depth), interaction (rendezvous, priced by the
communication) or reflection (quotation). The grades are separated in practice
by a choice of namespace, not by a modality: the reason a threat marker in the
noughts listing lives on its own name family is that were a line to signal a
threat on a square's move channel, the signal would rendezvous with the
square's receipt and the line would have \emph{played the move}. Perception must
not be able to act.

\paragraph{The generated logic.} One primitive modality $\Kmod{j}$ per rewrite
rule and choice of redex position, with $j$ ranging over redex positions;
structural connectives including separating conjunction; the name predicate
$\qt\varphi$; graded separation over a \emph{described namespace},
$\varphi \gsep{k}{N} \psi$; greatest fixed point $\nu$. There is no hidden-name
quantifier and no need of one, because \texttt{new} is syntactic sugar: a name
is a quoted process, no process contains the name that codes it, so a name
fresh in $P$ is manufactured by quoting any term properly containing $P$.
Consequently every family of names used below is \emph{described} by a
predicate rather than enumerated, which is what makes graded separation over it
available.

\paragraph{The assay.} A hypothesis is not consulted and then acted on. It sits
in guard position in a \texttt{where} clause, so acting \emph{is} testing and
confirmation \emph{is} the rendezvous. The complementary guard pair yields three
outcomes --- confirm, refute, and a third that is empty, in which the budget was
spent and nothing was learned. The third is ignorance, not refutation, and it is
indistinguishable from the environment being faulty. That identification of the
two bottoms is imported wholesale.

\paragraph{Conservation and predation.} Nothing mints. The arena routes. On a
win the arena hands the winner the loser's metabolic name, so $\sigma + \kappa =
\Theta$ throughout. An exogenous reward would need a referee with a mint and a
referee with a mint destroys conservation.

\paragraph{Read versus drive.} There are two routes to any predicate. Read it,
at perception prices, if the environment has already computed it into legible
structure; or drive it, priced by rendezvous. This is a claim about
environments, not about learners, and it is statable only because the framework
prices the grades separately and both parties have stacks. The
noughts-and-crosses measurement found that the three non-modal rungs of its
ladder delivered $+0.444$ of a total $+0.481$ gain --- $92.3\%$ of the yield ---
for $87.9$ of $255.9$ metered units, that is $34.3\%$ of the cost.

\subsection{From the consensus and quorum line}

Coalition logic in Gabbay's presentation, augmented with the behavioural
connective of context-labelled modalities, with the graded reading developed in
the quorum note: the coalition modality over a namespace $N$ valued in a
resolution semiring $V$,
\[
  \Emod{N}{V}\,\psi \;=\; \bigoplus_{O \subseteq N \text{ open}}
                          \bigotimes_{q \in O} \psi(q),
\]
where opens are described sub-namespaces via the name predicate. Three results
are imported. First, opens are not enumerated but described, so
$\varphi \gsep{k}{N} \psi$ applies to coalitions. Second, over $\Rnn$ with
$\otimes = \times$ a large coalition decays toward zero, so the correct quorum
semantics is $\otimes = \min$, a coalition being as strong as its weakest
member. Third, over the tropical and Viterbi algebras the attainment proposition
holds on the nose, so the coalition modality is \emph{attained} and the reader
gets a witness coalition rather than a scalar.

The two-quorum theorem of the swarm note is imported in statement and
re-proved in \S\ref{sec:coord} in a form specific to this setting, where it is
much stronger.

\subsection{From the platform}

Cost-accounted \rhoc\ as it stands on \texttt{feature/cost-accounted-rho}. Two
facts about the native integration bear on the spec. Metering is per-COMM at the
reducer, keyed by the deploy envelope's signature, and a signed term is
\emph{recognised} rather than lowered to a parking gate, so the normalised
\texttt{Par} has the same COMM count as the unsigned program and there is no
double-metering. In-program parking does not occur natively; an unfunded deploy
is rejected at the acceptance gate rather than parked mid-program. The surface
forms are unchanged: \texttt{s :: ()} mints, \texttt{\{\% P \%\}[s]} gates,
\texttt{(*)} is joint authority, \texttt{::} stacks a budget, \texttt{-o}
delegates, \texttt{\# P} signs a section.

The live guard fragment is arithmetic and pattern only. Structural and modal
guards are marked \texttt{//! ideal} in every listing in the repository, meaning
they specify a check rather than perform one. Appendix~\ref{app:rho} follows
that convention exactly.

\subsection{What is assumed to exist at development time}
\label{sec:assumed}

The instruction under which this note is written is that an implementation of
coalition logic together with context-labelled modalities will be available. We
therefore do not design around its absence, but we do state precisely what is
being relied on, because a spec that assumes a component should say what the
component must do. These are the requirements; they are collected again as an
interface in \S\ref{sec:checker}.

\begin{requirement}[Satisfaction, not synthesis]
\label{req:direction}
The checker must answer $P \models \varphi$. The consensus note's term
assignment runs the other way, $\vdash \varphi \rhd P$, proofs as processes; the
society needs the model-checking direction, because satisfaction is the guard
side and a guard side is what a \texttt{where} clause is. Coalition logic here
gives a piece a way to \emph{ask} whether a coalition has formed, not a way to
build one. This is the reason the whole synthesis-side gap of the consensus note
--- its missing proof system --- is not on this project's critical path.
\end{requirement}

\begin{requirement}[Three-valued verdicts with a budget bottom]
Verdicts are Kleene $\{f < b < t\}$ with $b$ carrying both readings: the
participant is faulty or stuck, and the observer could not afford to decide.
The caller must not be able to distinguish them, and the spec below depends on
that indistinguishability being visible in the measurements.
\end{requirement}

\begin{requirement}[Semiring-parametric grading with attainment]
The value algebra is a parameter. At minimum $\mathbb{B}$, $\Rnn$, tropical
$(\min,+)$ and Viterbi $(\max,\times)$ must be available. For the tropical and
Viterbi instances the checker must return a \emph{witness}: the coalition
attaining the $\oplus$, not merely its value.
\end{requirement}

\begin{requirement}[Indexed context-labelled modalities]
$\Kmod{j}$ must be indexed by redex position, and the quantifiers $\exists j$
and $\forall j$ must range over redex positions of the presented term. An
action-labelled logic will not do: every claim in the noughts world was the same
action, and in chess every quiet move is the same rewrite rule. The distinctions
that matter are positional.
\end{requirement}

\begin{requirement}[Settled modality]
A derived modality $\KKmod{j}$ meaning ``the rewrite at $j$ fires, the board's
deposit computation settles, and then $\varphi$ holds'' must be available, with
the settling shown terminating and confluent. In the noughts world this was
Lemma~8; \S\ref{sec:board} restates the obligation for chess, where it is
harder because ray scans are recursive.
\end{requirement}

\begin{requirement}[Costed queries]
Every query returns a metered cost alongside its verdict, charged to the
caller's stack, so that the ladder of \S\ref{sec:ladder} can be priced and the
measurements of \S\ref{sec:measure} can be taken. A checker that answers for
free would make the entire experiment vacuous.
\end{requirement}

\begin{requirement}[Namespace-indexed opens]
The coalition modality must accept its namespace as a name predicate, so that
``my two rooks'' and ``the pieces bearing on \texttt{f7}'' are expressible
without enumeration, and so that graded separation over the coalition namespace
is available.
\end{requirement}

\begin{requirement}[Caching by verdict]
Repeated queries of the same formula against an unchanged deposit must be
served from cache at the price of a lookup. Sixteen agents will ask overlapping
questions; without caching the read-versus-drive measurement of
\S\ref{sec:measure} measures the checker's implementation rather than the
architecture.
\end{requirement}

\section{Why chess, and why now}
\label{sec:why}

The programme's existing worked example is noughts-and-crosses. Every number in
it is exact over the $255{,}168$ plays, which is its strength, and it is solved
and tiny, which is its weakness. Four things are true of chess and false of
noughts, and each of them turns a stipulation of the earlier note into a
question with an answer.

\paragraph{Material exists, so eating is non-trivial.} In noughts nothing is
captured; predation happens only at the arena, between games, when a loser's
metabolic name is handed over. In chess predation happens \emph{during} the
game, at the granularity of a single ply, and the prey is a peer agent with its
own stack. Conservation has to hold across that, and the accounting becomes
interesting rather than ceremonial.

\paragraph{Coalitions are real, so coalition logic has content.} Chess strategy
is already coalition talk: the rook and bishop together force mate, these three
pieces hold the kingside, that knight is doing two jobs. In noughts a coalition
is a pair of threats and the vocabulary is overkill. In chess the coalition
modality has something to quantify over.

\paragraph{The position is legible, so read-versus-drive is a real choice.} The
attack map of a chess position is expensive, shared, and recomputable. Whether
the world deposits it or each agent drives it is a genuine design fork with a
measurable cost difference. In noughts the analogous object is eight lines and
the question does not bite.

\paragraph{The science cannot be closed, so the ecology does not die.} The
noughts note found that $\chi$, the exact characteristic formula, is reached at
nine unfoldings, and then found that this is fatal: a $\chi$-monoculture draws
every game, the harvest rate goes to zero, and the population starves on a
finite wall. A solved science is a dead ecology and imperfection is load-bearing.
In chess $\chi$ is not reachable at any budget the framework admits, so the
ecology cannot die that way. We record this as the reason the ecological half of
the programme becomes \emph{observable} here for the first time.

\begin{remark}[Chess keeps the ecology alive]
\label{rem:alive}
The noughts result that a solved science is a dead ecology was obtained by
exhibiting the solved case. It is a theorem about a closed science and its
practical content --- that real populations must be imperfect to eat --- could
not be observed in a world where perfection was affordable. Chess supplies a
world where it is not. This does not make the earlier result stronger; it makes
it testable in a regime where the alternative hypothesis, that the population
converges on a single good policy and stops, is actually available.
\end{remark}

Against those four, one thing is true of noughts and false of chess: exactness.
Nothing in the chess project will be exact over the game tree. Every measurement
in \S\ref{sec:measure} is a sampled estimate against a fixed reference opponent,
and the note should say so in its own results section rather than discovering it
in review.

\section{Two societies, one board}
\label{sec:two}

The proposal names two things as agents: the sixteen pieces and the board. It is
worth being explicit that these are two \emph{societies}, with different jobs,
different budgets, and different logical characters.

\begin{definition}[The world and the player]
\label{def:two}
The \emph{world} is the society of sixty-four square-agents together with the
ray, attack, pin and legality machinery they support, plus the arena and the
clock. The \emph{player} is the society of at most sixteen piece-agents of one
colour. The world holds no strategic stack and mints nothing; it deposits
structure and refuses illegal joins. The player holds the budget and makes the
decision.
\end{definition}

This is the same split the noughts listing already makes, between
\texttt{Board}, \texttt{Cells}, \texttt{Line} and \texttt{Scan} on one side
and \texttt{Choose}, \texttt{Move} and \texttt{Revise} on the other, and it is the split
that makes the question \emph{who is billed} askable. The noughts note observed
that putting threats in the term relocates the cost of the scan from the
scientist's budget to the world's, and that the interesting question is who pays.
In chess that question has a large answer, because the object in dispute is the
attack map.

\begin{remark}[The decomposition is a choice, and there are others]
\label{rem:decomp}
Agent-per-piece is not the only decomposition. Agent-per-square is the world.
One could also try agent-per-ray (the sliding lines as agents, which is close to
how the noughts \texttt{Line} contract behaves), or agent-per-plan, where an
agent is a coalition and pieces are its resources. Agent-per-piece is chosen
because it is the decomposition in which predation is definable: a piece can be
captured and a ray cannot. The alternatives are worth naming because
agent-per-plan is what a strong human player's introspection reports, and it is
the natural second experiment.
\end{remark}

\section{The board}
\label{sec:board}

\subsection{Why one name family per square is not enough}

The proposal's encoding is
\[
  \prod_{1 \leq i,j \leq 8} \qt{(i,j)}\,!\,(\,\cdot\,)
\]
with the payload $\Nil$ for unoccupied and $*x$ for the fresh channel minted to
interact with a particular piece. The occupancy half of this is right and the
manufactured-name half is right --- we return to both --- but the encoding
conflates two roles that the existing listings deliberately separate, and the
reason they separate them is stated in the noughts source: perception must not be
able to act, and that separation of access grades is enforced by choice of
namespace.

In noughts the separation is cheap because perception is short-range: eight
lines read three cells each. Chess has genuine long-range perception. A rook on
\texttt{a1} must know what stands on \texttt{a8} without any intention of going
there, and a bishop's diagonal passes over squares it will never occupy. If
reading occupancy means consuming and re-emitting $\qt{(i,j)}!(v)$, then every
perception is billed as an interaction, sixteen agents scanning concurrently
serialise on every square, and the read-versus-drive distinction --- which is one
of the two things this project exists to measure --- has been designed out of the
system before the first line of code.

\begin{definition}[Three name families per square]
\label{def:families}
Let $b$ be the board's channel and $s$ a square index. The board maintains:
\begin{description}[leftmargin=2.2em,style=nextline]
\item[the occupancy cell] $\mathit{occ}(s) = \qt{[*b,\texttt{"occ"},s]}$, a
resting send carrying either $\Nil$ or $*p$ where $p$ is the piece-agent's
channel. This is the only family a ply consumes.
\item[the attack family] $\mathit{atk}(s,c,p) =
\qt{[*b,\texttt{"atk"},s,c,p]}$, standing offers deposited by the settling
computation, one per attacker $p$ of colour $c$ bearing on $s$. Read-only in
practice: no move protocol receives on it.
\item[the constraint family] $\mathit{pin}(s)$, $\mathit{chk}(c)$,
$\mathit{ep}(s)$, $\mathit{cast}(c,\mathit{side})$, likewise deposited, carrying
the derived facts a legality guard consults.
\end{description}
All three are manufactured names in the sense of the freshness-by-quotation
proposition, not \texttt{new}-bound, so each family is describable by a name
predicate and graded separation over it is available.
\end{definition}

\begin{proposition}[Perception must not act]
\label{prop:percept}
If attack markers are deposited on the occupancy family, then a scan that
signals ``$p$ bears on $s$'' can rendezvous with a move protocol receiving on
$\mathit{occ}(s)$, and the scan performs a ply. Separating the families makes
that failure unrepresentable rather than merely unlikely.
\end{proposition}

\begin{proof}
Immediate, and it is the chess instance of the argument written into the noughts
source for threat offers. The content is that the separation is structural: with
distinct families there is no term in which a scan and a move protocol can
rendezvous, so the property does not depend on any discipline the scan follows.
\end{proof}

\subsection{Who is billed for the attack map}
\label{sec:billed}

The attack map is the expensive object in chess and it is shared by every agent
in the player. There are two architectures.

\begin{description}[leftmargin=1.6em,style=nextline]
\item[World-billed.] After each ply the board runs a settling computation that
walks every ray and deposits $\mathit{atk}$ markers. The player's agents then
read markers at perception prices. One computation, sixteen readers.
\item[Player-billed.] The board deposits only occupancy; each agent that needs
to know whether a square is attacked drives the query itself, paying a
rendezvous per square walked. Sixteen computations of one fact.
\end{description}

The framework's own thesis says the first should win, and by a factor near the
population size, because a predicate the environment has computed into legible
structure is read at formula-depth prices rather than driven at rendezvous
prices. But the framework does not say this is free: the world's settling
computation is metered too, and the noughts note's remark on the point is
precisely that the cost is \emph{relocated}, not removed, and that the
interesting question the framework can pose --- because both parties have stacks
--- is who is billed. We therefore specify both and measure the difference
(measurement M7, \S\ref{sec:measure}). This is the cheapest genuinely
informative experiment in the whole project.

\begin{remark}[Settling in chess is harder than in noughts]
\label{rem:settle}
The noughts settling lemma had an easy proof: each of eight lines fires once, no
line's output enables another because offers go to names no line reads, so the
computation terminates and is confluent. The chess settling walks rays, and a
ray's walk is recursive in the board's occupancy: a slider's ray terminates at
the first occupied square. Termination is still clear --- each ray step consumes
one square of a strictly decreasing remaining ray --- and confluence still holds
for the same reason, that markers are deposited on a family nothing in the
settling reads. But the obligation must be discharged explicitly, because
$\KKmod{j}$ depends on it and so does every modal rung of the ladder.
\end{remark}

\subsection{Legality is a universal, and it belongs to the world}
\label{sec:legality}

Check, pin, and discovered check are constraints on the legality of \emph{every}
piece's move, and they are global: whether the bishop on \texttt{c4} may move
depends on whether doing so exposes the king. No individual piece can establish
this locally and no quorum of pieces can establish it either, because an illegal
move is not a move and a vote does not make it one.

\begin{principle}[The legality universal is the world's obligation]
\label{prin:legality}
The one unavoidable $\forall$ in chess is placed on the world's budget, not the
player's. The world refuses illegal joins, the way the arena in the noughts
world holds no stack and only routes. The consequence is that the player's own
hypothesis language stays existential and finitely confirmable, which is the
regime in which the asymmetry proposition says witnesses are cheap.
\end{principle}

There is a second consequence, and it is the one that makes this satisfying
rather than merely convenient. If the world deposits $\mathit{pin}(s)$ markers
--- computed by walking outward from each king through the first own piece to
the first enemy slider --- then pin-awareness becomes perception. A piece does
not reason about whether it is pinned; it reads whether a marker stands on its
square. Likewise check is a marker, not a search. The single most classically
``knowledge-heavy'' part of chess legality becomes deposited structure, and the
read-versus-drive thesis gets its best instance.

\begin{remark}[What the world does not deposit]
\label{rem:notdeposit}
The world deposits facts about the \emph{current} position: occupancy, attacks,
pins, checks, en-passant and castling rights. It does not deposit facts about
\emph{consequences} of moves, because those are the player's hypotheses and
pricing them to the world would smuggle the evaluation function back into the
architecture through the floor. The line between the two is exactly the line
between the non-modal and the modal rungs of the ladder, and it should be
guarded in review: any proposal to have the board deposit ``this move loses a
piece'' is a proposal to write a conventional engine and call it an ecology.
\end{remark}

\section{A ply is a join}
\label{sec:join}

\subsection{What the noughts construction got, and why chess does not inherit it}

The second draft of the game note was a correction: the context $K$ had never
been instantiated, and instantiating it as a \emph{redex} context was what made
$\Kmod{s}$ a generated primitive rather than a derived abbreviation. The payoff
was a bijection. An empty square is a receipt offering to be claimed; a claim is
one join consuming the mover's send and the stale mark; after it fires no
receipt remains. So redex positions correspond to empty squares correspond to
legal moves, and the branching factor of $\exists s$ \emph{is} the legal move
count.

None of that transfers to chess for free. A chess move touches two squares at
minimum, vacating one and occupying another. A capture additionally ends a
sibling agent's computational life. Castling touches four squares and two
pieces. En passant has a destination distinct from the square of the captured
pawn. Promotion replaces a piece's identity. If a ply is several rewrites then
$\Kmod{j}$ indexes something finer than a move, the bijection fails, the modal
rungs of the ladder index the wrong thing, and --- worse operationally --- two
plies can interleave in the middle of each other.

\subsection{The repair}

The repair is available and it is already demonstrated in the platform. The
cost-accounting demonstration uses an $N$-ary atomic join for a four-cell swap
that conserves money and goods all-or-nothing in every interleaving, with the
guard deciding before any leg fires. A chess ply is that shape exactly.

\begin{definition}[A ply]
\label{def:ply}
A ply is a single $N$-ary join over the occupancy cells it touches, guarded by a
\texttt{where} clause consulting the constraint family, which re-emits every
consumed cell with its updated payload. Formally, for a move $m$ touching
squares $s_1,\dots,s_n$,
\[
  \texttt{for}\;(\,\qt{v_1} \leftarrow \mathit{occ}(s_1)\;\&\;\cdots\;\&\;
  \qt{v_n} \leftarrow \mathit{occ}(s_n)\;\texttt{where}\;
  \mathsf{Legal}(m,\vec v)\,)\;\{\,\cdots\,\}
\]
where the body re-emits $\mathit{occ}(s_1),\dots,\mathit{occ}(s_n)$ with the
post-move payloads, terminates any captured piece-agent, and triggers the
settling.
\end{definition}

\begin{proposition}[Move taxonomy is join arity]
\label{prop:arity}
Under Definition~\ref{def:ply} the classification of chess moves by arity and
payload transformation is exactly the classical move taxonomy:
\begin{center}
\begin{tabular}{@{}lll@{}}
\toprule
Move kind & Arity & Payload transformation \\
\midrule
Quiet move          & 2 & $(*p,\Nil) \mapsto (\Nil,*p)$ \\
Capture             & 2 & $(*p,*q) \mapsto (\Nil,*p)$, $q$ terminated \\
En passant          & 3 & $(*p,\Nil,*q) \mapsto (\Nil,*p,\Nil)$, $q$ terminated \\
Castling            & 4 & $(*k,\Nil,\Nil,*r) \mapsto (\Nil,*k,*r,\Nil)$ \\
Promotion           & 2 & $(*p,\cdot) \mapsto (\Nil,*p')$, $p$ retired, $p'$ born \\
\bottomrule
\end{tabular}
\end{center}
\end{proposition}

We state this as a proposition rather than a table because it has a consequence.

\begin{corollary}[$\Kmod{j}$ survives]
\label{cor:Ksurvives}
With plies as joins, redex positions of the ply rule correspond bijectively to
legal moves, so $\Kmod{j}$ remains a generated modality with $j$ ranging over
moves, $\exists j$ has the legal move count as branching factor, and every modal
rung of the ladder indexes what it is supposed to index. The whole ladder
apparatus of the game note transfers.
\end{corollary}

\begin{remark}[The alternative, and why not]
\label{rem:alt}
One could let a ply be several rewrites and recover a composite via the settled
modality, as the noughts world does for the board's own settling. The cost is
that a ply is then no longer atomic in the term, so two agents' plies can
interleave, and the guard that makes a move legal would have to be re-checked
against a position that may have changed under it. Chess has a turn structure
and it should be represented by atomicity rather than by convention. The join
route also costs nothing: the platform has $N$-ary joins and the demonstration
of a conserving four-cell swap already exists.
\end{remark}

\begin{remark}[Promotion is birth and it is not free]
\label{rem:promo}
Promotion is the only move that creates an agent. Under conservation it cannot
create a \emph{budget}: the promoted agent's stack must come from somewhere, and
the natural source is the pawn's own stack, transferred whole. That makes
promotion a change of policy and capability at constant budget, which is the
right reading, and it means the classical observation that a promoted queen is
worth a queen is not automatic here --- it is a prediction, since the promoted
agent inherits a pawn's accumulated stack rather than a queen's.
\end{remark}

\section{The coordination problem}
\label{sec:coord}

\subsection{The shape of the problem}

Sixteen agents each want to move; exactly one move commits. This is not a
technical inconvenience, it is the whole architectural content of the proposal,
and the programme already contains the analysis of it.

\begin{observation}[The K/J split]
\label{obs:kj}
A graded clause is a map from candidates to values: a \emph{weighting} of the
fibre, not a \emph{section} of it. In the selection-monad reading, a clause is a
quantifier $(X \to R) \to R$ and a resolver is a selection function
$(X \to R) \to X$. The swarm exhibits the split as different behaviours
performed by different bees: dance evaluation is the clause, piping and takeoff
are the section, and no bee does both.

The piece society must exhibit the same split. A piece \emph{evaluates}
candidate moves; something else \emph{commits}. A design in which each piece
both scores and plays is a design with sixteen resolvers and no game.
\end{observation}

\subsection{Quorum, not consensus, and the argument is not a preference}

The obvious architecture is a council: all sixteen pieces report scores, the
scores are combined, the argmax commits. That architecture has a chairman, and a
chairman is a global evaluator wearing a hat. The programme's altitude argument
rules it out, and it does so by derivation rather than by taste.

\begin{proposition}[The mechanism is forced]
\label{prop:forced}
Consensus at the council is a global $\forall$ over all sixteen advocates with
no presented witness: finitely refutable, not finitely confirmable. Quorum at a
candidate move is a local $\oplus\otimes$ over one open. No participant can
evaluate the global universal within a bounded budget, so a society of
budget-bounded participants must use the local formula.
\end{proposition}

This is the quorum note's derivation of Seeley's empirical finding, restated for
a designed rather than an evolved population. In the swarm it explains why
quorum is sensed at the site rather than at the cluster. Here it says: pieces
advocate for candidate \emph{moves}, and a candidate that accumulates enough
advocacy commits. Nobody counts the whole distribution.

\subsection{The advocacy namespace, and why a candidate move is a name}

\begin{definition}[Advocacy]
\label{def:advocacy}
A candidate move $m$ is a term; by reflection $\qt{m}$ is a name. The advocacy
channel for $m$ is $\mathit{adv}(m) = \qt{[*\mathit{pl},\texttt{"adv"},\qt{m}]}$
where $\mathit{pl}$ is the player's channel. A piece advocates for $m$ by
sending its advocacy token on $\mathit{adv}(m)$. An \emph{open} in the coalition
sense is the set of pieces that have advocated for a given $m$ during the
current round; it is described by a name predicate, not enumerated.
\end{definition}

This is the herd note's corollary landing exactly: opens are described
namespaces via $\qt\varphi$. It is also, we think, the most pleasing single
consequence of working in a reflective calculus. In the swarm an open is a set
of bees and the modeller has to say so. Here the site a piece advocates for
\emph{is} a name, manufactured by quoting the move, and the advocacy set is a
sub-namespace of the player's namespace. Nothing has to be introduced to make
coalition logic apply; the calculus already provides the points.

\subsection{One token per piece per round: the society is a polity}

\begin{definition}[The franchise]
\label{def:franchise}
At the start of each advocacy round the player issues to each surviving piece
exactly one \emph{linear} advocacy token, of signature $\mathsf{V}$, into a
located stack whose channel is manufactured from the piece's identity. A piece
spends its token by advocating, or by inhibiting, or not at all. Tokens do not
carry over between rounds.
\end{definition}

This is the polity listing's pre-election phase transposed without change: a
persistent claim on the roll metered to one linear token per epoch, a source
keyed to identity and not a mint. The reason to adopt it is not tidiness. It is
that it makes the following theorem an arithmetic fact about sends rather than a
statistical regularity.

\subsection{Quorum is a join arity}

\begin{definition}[The commit gate]
\label{def:gate}
The commit gate for candidate $m$ at threshold $k$ is the $k$-ary join
\[
  \texttt{for}\;(\,\qt{a_1} \leftarrow \mathit{adv}(m)\;\&\;\cdots\;\&\;
  \qt{a_k} \leftarrow \mathit{adv}(m)\,)\;\{\;\mathsf{Commit}(m)\;\}
\]
which fires exactly when $k$ distinct advocacy sends stand on $\mathit{adv}(m)$,
consuming all of them.
\end{definition}

\begin{theorem}[Two quorums need $2k$ tokens]
\label{thm:2k}
Let $S$ be the number of surviving pieces, each holding exactly one linear
advocacy token per round, and let $k$ be the commit threshold. Then two distinct
candidate moves cannot both fire their commit gates in the same round unless
$S \geq 2k$. In particular with $S = 16$ and $k \geq 9$ a split is impossible.
\end{theorem}

\begin{proof}
Each commit gate consumes $k$ sends on its own channel. Tokens are linear, so a
token spent on $\mathit{adv}(m)$ is not available on $\mathit{adv}(m')$ for
$m' \neq m$, and each piece holds one. Two firings therefore consume $2k$
distinct tokens drawn from a pool of $S$.
\end{proof}

\begin{remark}[What this theorem is worth]
\label{rem:2kworth}
In the swarm the corresponding statement is a structural impossibility result
that a statistical model cannot state, and it was confirmed by simulation:
splits identically zero at $S = 12,16,20,24,28,29$ and first appearing at
$S = 30 = 2k$ with $k = 15$. Here it needs no simulation, because linearity of
the token and arity of the join do the counting. That is the difference between
importing a result and re-deriving it in a calculus that already had the right
primitives. It also means the designer has a free safety property: choose
$k \geq \lceil S/2 \rceil$ and the society \emph{cannot} commit two moves, for
any policies whatsoever, with no cross-inhibition machinery at all.
\end{remark}

\subsection{The other failure mode, and what cross-inhibition is for}

Safety comes free; liveness does not. With $k \geq 9$ and advocacy split eight
against eight, no gate fires and the round deadlocks. The swarm literature's
division of labour applies exactly: below $2k$ cross-inhibition prevents
\emph{deadlock}; above $2k$ it prevents \emph{split}. Since the design fixes
$k \geq S/2$ and buys no-split structurally, cross-inhibition in this setting
has only the deadlock job, and it competes with a cheaper mechanism.

\begin{definition}[Rounds, and the deposited distribution]
\label{def:rounds}
A turn consists of up to $R$ advocacy rounds. At the end of a round in which no
gate fired, unconsumed advocacy sends remain standing on their
$\mathit{adv}(m)$ channels and are \emph{readable}. Fresh tokens are issued. A
piece may therefore re-advocate having perceived, at perception prices, where
the rest of the society stood.
\end{definition}

This is the bees reading dance strength at the cluster: cheap perceptual price
for the reversible decision (which candidate to back), expensive interaction
price for the irreversible one (the commit). We take that as the default and
treat cross-inhibition as a labelled variant, for two reasons. It is one more
mechanism, and the swarm measurement recorded an unexpected cost --- with cross
inhibition on, the swarm chose the better site slightly \emph{less} often,
because inhibition suppresses correct advocacy too.

\begin{definition}[The resolver of last resort]
\label{def:lastresort}
If $R$ rounds elapse with no commit, the clock commits the candidate carrying
the most standing advocacy, ties broken by a deterministic order on $\qt{m}$.
\end{definition}

This is the one centralised element in the design and the note should not hide
it. It is held by the clock, which is world and not player; it has no evaluation
function and no preferences; and its firing rate is a measurement
(M4, \S\ref{sec:measure}), because a society that reaches last resort often is a
society whose quorum parameters are wrong.

\subsection{The population shrinks, and it crosses the boundary during play}

\begin{proposition}[The endgame crossing]
\label{prop:crossing}
$S$ is monotonically non-increasing over a game, from $16$ toward as few as one
piece besides the king. If $k$ is held fixed as an absolute number, then
$S < k$ eventually holds and no gate can ever fire; if $k$ is held fixed as a
fraction $\alpha S$ with $\alpha > 1/2$, splits remain impossible throughout but
the absolute number of advocates required falls. Either way the society crosses
a regime boundary at a computable ply.
\end{proposition}

\begin{remark}[A prediction the bees could not supply]
\label{rem:predict}
The swarm's $S$ is fixed within an episode, so the two-quorum theorem there
partitions \emph{episodes}. Chess partitions the \emph{inside} of a single
episode: the same society passes through the split-possible regime and into the
split-impossible one as material comes off. The framework therefore predicts
that the character of the decision procedure must change at a computable point,
with competing plans available in the middlegame and structurally unavailable in
the endgame. That this matches ordinary chess intuition --- middlegames are
about competing plans, endgames about executing one --- is not evidence, but it
is the kind of coincidence that makes a measurement worth taking rather than a
post-hoc story worth telling. M5 in \S\ref{sec:measure} is that measurement, and
it is the one we would report first.
\end{remark}

\section{Coalition logic as the society's strategic vocabulary}
\label{sec:coalition}

\subsection{The vocabulary is already chess's own}

Chess commentary is coalition talk. The rook and bishop together force mate.
These three pieces hold the kingside. That knight is overloaded, doing two jobs
at once. That bishop is bad, belonging to no coalition. Coalition logic's
$\langle E \rangle\varphi$, ``coalition $E$ can force $\varphi$'', is not a
formalism imported and then applied; it is the native register of the domain.

This matters for the project because it means the society's hypotheses can be
\emph{about} coalitions without anyone having designed a coalition-formation
mechanism. Pieces advocate; advocacy sets are opens; opens are coalitions; and
the logic evaluates formulae over them.

\subsection{Which algebra, and why the answer is not $\Rnn$}

\begin{proposition}[Large coalitions must not decay]
\label{prop:decay}
Over $\Rnn$ with $\otimes = \times$ and values in $[0,1]$, a coalition's value
$\bigotimes_{q \in O}\psi(q)$ is non-increasing in $|O|$ and strictly decreasing
whenever any member's value is below one. Hence larger quorums score weaker,
which inverts the intended semantics.
\end{proposition}

\begin{definition}[Quorum algebras]
\label{def:quorum}
A quorum algebra takes $\otimes = \min$: a coalition is as strong as its weakest
member. Over the tropical $(\min,+)$ and Viterbi $(\max,\times)$ algebras the
attainment proposition holds on the nose, so $\bigoplus$ is realised by a
definable selection function.
\end{definition}

We note, because it is a pleasant piece of evidence rather than an argument, that
$\otimes = \min$ is independently the correct chess semantics. A combination
fails at its weakest participant. If three pieces bear on a mating square and one
of them is pinned, the combination is worth what the pinned piece is worth, which
is nothing; the other two do not compensate. Overloading is the same shape from
the other side: a piece that belongs to two coalitions caps both at its own
availability. The algebra the framework forces on general grounds is the algebra
the domain would have chosen.

\begin{corollary}[The player gets a witness]
\label{cor:witness}
Over the attained algebras the coalition modality returns not a scalar but the
coalition realising it. The society therefore commits to a move \emph{together
with the coalition that justifies it}.
\end{corollary}

\begin{remark}[Explanation is structural, not decorative]
\label{rem:explain}
Corollary~\ref{cor:witness} is the closest thing this project has to a
by-product worth having on its own. An explanation of the form ``this move,
because these pieces, at this strength'' is not generated after the fact by a
separate narration module reading the engine's numbers. It is the witness the
selection function returned. Whatever else is true of the resulting player, its
explanations cannot be confabulated, because there is no separate faculty
available to confabulate with.
\end{remark}

\subsection{The alternating rung, and the question this project exists to answer}
\label{sec:altrung}

Coalition logic's ``can force'' quantifies over the opponent's replies. It is
therefore an alternating formula, a box under a diamond, and the game note's
measurement of the corresponding rung was unkind: rung~4, the one genuinely
alternating rung of the noughts ladder, bought $+0.005$ for $111$ metered units
and was the worst purchase on the ladder, while the three non-modal rungs
delivered $92.3\%$ of the yield for $34.3\%$ of the cost.

We do not think this settles anything, and the reason is the reason for the
whole project.

\begin{observation}[The yield result is a measurement on one game]
\label{obs:onegame}
Noughts-and-crosses has $5478$ reachable positions, no material, no long-range
interaction, and a branching factor that starts at nine and falls. Its non-modal
predicates --- a completing threat, a fork, centrality --- are close to
sufficient because the game is close to trivial. Whether the finding generalises
is not a question the noughts data can answer, because the finding could be a
fact about games or a fact about the smallness of that game.
\end{observation}

\begin{principle}[Chess is the discriminating instance]
\label{prin:discriminate}
The measurement that would most change the programme's view of itself is the
ladder measurement (M1, M2) repeated on chess. If the non-modal rungs again
deliver most of the yield for a third of the cost, the read-versus-drive thesis
acquires its second and much harder instance and becomes a claim worth defending
in general. If the modal and coalition rungs dominate in chess, then the noughts
result is a result about noughts, and the programme has learned something it
could not have learned from a game it could solve.
\end{principle}

This is the strongest strategic argument for building the thing, and it should
appear in the note's own abstract, which it does.

\subsection{Graded separation and what a fork is in chess}

The noughts fork was $\mathsf{Threat}(p) \gsep{0}{N} \mathsf{Threat}(p)$: two
threats sharing no square, graded at separation zero over the square-indexed
namespace. Requiring it is what forced the board encoding, since two threats
with a common gap must be excluded and therefore squares rather than lines must
be the parands.

The chess transposition is direct and the namespace is the attack family.

\begin{definition}[Fork, chess instance]
\label{def:fork}
Let $N_{\mathit{atk}}$ be the namespace described by the name predicate
\[
  \qt{[*b,\texttt{"atk"},\cdot,c,\cdot]},
\]
and write $\mathsf{Bears}(p,s)$ for the structural predicate asserting a marker
deposited by $p$ on $s$. Then
\[
  \mathsf{Fork}(p) \;:=\;
  \big(\exists s.\,\mathsf{Bears}(p,s) \wedge \mathsf{Valuable}(s)\big)
  \gsep{0}{N_{\mathit{atk}}}
  \big(\exists s'.\,\mathsf{Bears}(p,s') \wedge \mathsf{Valuable}(s')\big).
\]
\end{definition}

The same argument that forced squares as parands in noughts forces
\emph{squares} here too, not pieces: a knight bearing twice on one square is not
a fork, and separation over a piece-indexed namespace would not exclude it.

\subsection{The coalition rung}

\begin{definition}[The coalition rung]
\label{def:coalrung}
For a target square $s$ and the player's piece namespace $N_{\mathit{pl}}$,
\[
  \mathsf{Bearing}(s) \;:=\;
  \Emod{N_{\mathit{pl}}}{V}\; \lambda p.\;\mathsf{Bears}(p,s)
  \otimes \mathsf{Available}(p),
\]
where $\mathsf{Available}(p)$ is zero if $p$ is pinned or already committed
elsewhere. Over a quorum algebra this evaluates to the strength of the strongest
available coalition bearing on $s$, attained at a witness.
\end{definition}

$\mathsf{Available}$ is where pins enter, and it is worth observing that it
enters as a \emph{factor} in the $\otimes$ rather than as a side condition. A
pinned piece contributes $\bot$ to any coalition containing it, and under
$\otimes = \min$ that zeroes the coalition. This is the correct behaviour and it
was not designed in; it follows from the choice of algebra forced in
\S\ref{sec:coalition}.

\begin{remark}[Brittleness, and the two dials]
\label{rem:brittle}
The herd note recorded a negative result that applies here without change: $\min$
over a large coalition is zeroed by one dissenter, so a crisp rung and a graded
rung can return identical rows, and grading pays only if coalition \emph{size}
is relaxed as well. The chess instance is that $\mathsf{Bearing}$ over a fixed
large coalition is brittle, and the design must therefore relax over subsets ---
which it does, since $\bigoplus$ ranges over opens, and the witness returned is
the best available subset rather than the whole. We flag it because it was a
surprise in the earlier work and would be a surprise again.
\end{remark}

\section{Material as metabolism}
\label{sec:material}

The noughts arena routes and never mints: on a win the arena hands the winner the
loser's metabolic name. In chess, capture is predation at ply granularity, and
nothing needs reinterpretation for it to be so.

\begin{definition}[Capture is harvest]
\label{def:harvest}
When a ply's join terminates a piece-agent $q$, the capturing agent $p$ receives
$q$'s metabolic name; $q$'s remaining stack is thereby routed to $p$. No tokens
are minted and none destroyed, so $\sigma + \kappa = \Theta$ holds across the
ply.
\end{definition}

\begin{proposition}[Material advantage is budget advantage]
\label{prop:matbudget}
Under Definition~\ref{def:harvest}, a side with more surviving material has, all
else equal, both more advocates and more aggregate stack, hence can afford more
perception per turn and more rounds of advocacy. The conventional evaluation
term ``material'' becomes a statement about the society's computational capacity
rather than a number added to a score.
\end{proposition}

This is the most interesting thing in the design and also the least certain.

\begin{conjecture}[Emergent piece values]
\label{conj:values}
Let $\pi(t)$ be the average lifetime harvest --- tokens acquired by capture,
plus tokens the agent's survival makes available to the society --- of an agent
of piece type $t$ over a population of games. Then the ratios
$\pi(\text{pawn}) : \pi(\text{knight}) : \pi(\text{bishop}) :
\pi(\text{rook}) : \pi(\text{queen})$ approximate the classical
$1 : 3 : 3 : 5 : 9$.
\end{conjecture}

We state it as a conjecture and not as a design goal, and the difference matters.
Nothing in the architecture hands a piece its value. A queen is an agent whose
move-generation contract reaches further, which means it can advocate for more
candidates, bear on more squares, and belong to more coalitions; and under
$\otimes = \min$ belonging to more coalitions is worth something only if the
agent is available. If the classical ratios fall out, the framework has derived
one of the few quantitative facts about chess that has no derivation. If they do
not, the discrepancy is data about which of the framework's pricing choices is
wrong, and it is much more informative than agreement would be.

\begin{remark}[Two more in the same currency]
\label{rem:currency}
Offered as shapes rather than claims, because they are the sort of thing the
design makes sayable and it would be a waste not to say them. A \emph{pinned}
piece is one whose budget is committed but unspendable, which is why
$\mathsf{Available}$ zeroes it. \emph{Zugzwang} is being obliged to spend a token
when every available expenditure is loss-making, which in this setting is not a
metaphor: the franchise is issued per round and a piece that must advocate for
something is a piece that cannot decline the market.
\end{remark}

\section{Where the policy lives}
\label{sec:policy}

\subsection{The extremes are not a dial}

A piece's policy is the thing that maps the position to an advocacy. The
question is where that map lives and whether it is priced. The
sixteen-model-sessions extreme and the all-\flang\ extreme are two answers,
and they are not near-neighbours on a continuum.

A model session is opaque. Its choice function is not a clause, so it cannot sit
in guard position and the assay construction does not apply. It is not priced,
so cost accounting measures API calls rather than computation, and every
measurement in \S\ref{sec:measure} loses its denominator. It is not observable,
so the coalition formulae become commentary on a decision they did not make, and
Remark~\ref{rem:explain} is false --- a separate narrating faculty is exactly
what a model wrapper reintroduces. At that extreme every structural result in
this note evaporates and what remains is a multi-agent chess demonstration, of
which there are many. Sixteen sessions per advocacy round is also slow enough
that the round structure of \S\ref{def:rounds} becomes unusable.

\subsection{The aperture, which is the programme's own construction}

None of that makes model-backed pieces incoherent, and the reason is that the
choice-formulae note already worked out how strength enters a shard from
outside. Its Route~E is the \emph{aperture}: a channel fed by the environment
rather than by the term. Two things are said about apertures there. They are
where choice of full strength can enter, and they are syntactically conspicuous,
so a static analysis of aperture channels bounds a whole shard. Under the
stratified \texttt{where} discipline, subject reduction holds and full-strength
choice occurs exactly at declared apertures while the interior's choice type
stays invariant.

A model session is an aperture, precisely and not by analogy. So a society with
model-backed pieces is coherent in a stratified shard, and building it would
\emph{demonstrate} the stratification result on an artefact people care about
rather than on a counterexample two lines long.

\begin{principle}[Relocate, do not compromise]
\label{prin:relocate}
The recommendation is not a midpoint between the extremes. It is that the model
goes outside the ply loop by default, and inside only as a declared aperture in
a labelled variant.
\end{principle}

Outside the ply loop there are two jobs a model is actually good at, and both
are already roles in the existing design. The first is authoring and mutating
piece policies between games, which is the \texttt{Revise} and \texttt{Crossover}
and \texttt{Beget} role of the noughts listing, and which lands squarely on the
reproduction note's crossover over quoted code --- a model is a mutation operator
on programs, which is a thing models do well. The second is reading the witness
coalition out loud, which is post-hoc narration of a decision that was genuinely
made elsewhere and is therefore honest.

\begin{remark}[What the aperture variant would measure]
\label{rem:apvar}
Worth building second, not first, and worth building for a specific reason: it
is the only configuration in which one can measure what a model's judgement is
worth \emph{in the framework's own currency}. Run the same society with $j$ of
sixteen pieces aperture-backed and $16-j$ \flang-backed, and the yield-per-unit
metric prices the model's contribution against the ladder rungs. Nothing else in
the current AI-evaluation landscape prices a model's judgement against a
formally specified alternative on the same budget.
\end{remark}

\section{The implementation spec}
\label{sec:spec}

This section is written to be handed to a developer. It fixes names, protocols,
prices, interfaces, delivery phases and acceptance criteria. Where a decision is
genuinely open it is marked \textsc{open} rather than resolved silently.

\subsection{Module inventory}
\label{sec:modules}

\begin{center}
\begin{tabular}{@{}llp{6.6cm}@{}}
\toprule
Module & Society & Responsibility \\
\midrule
\texttt{Board}    & world  & occupancy cells; the ply join; refusal of illegal joins \\
\texttt{Settle}   & world  & ray walks; deposit of \texttt{atk}, \texttt{pin}, \texttt{chk}, \texttt{ep}, \texttt{cast} \\
\texttt{Arena}    & world  & game lifecycle; harvest routing; adjudication \\
\texttt{Clock}    & world  & turn and round structure; resolver of last resort \\
\texttt{Registry} & world  & the roll: which piece channels are live this turn \\
\texttt{Franchise}& player & issue of one linear advocacy token per live piece per round \\
\texttt{Piece}    & player & per-agent policy: generate candidates, evaluate, advocate \\
\texttt{Gate}     & player & the $k$-ary commit join, one instance per live candidate \\
\texttt{Ladder}   & player & the formula catalogue; rung held per agent \\
\texttt{Revise}   & player & per-agent rung revision from assay outcomes \\
\texttt{Metab}    & both   & \texttt{Live}: reclaim, debit, reissue; death on insufficiency \\
\texttt{Check}    & world  & the coalition-logic and modal checker service (\S\ref{sec:checker}) \\
\bottomrule
\end{tabular}
\end{center}

\texttt{Check} is placed in the world deliberately. It is an instrument the
player \emph{holds}, and its queries are charged to the caller, but it is not
part of any agent's term; otherwise sixteen agents would carry sixteen copies of
the checker and the caching requirement could not be met.

\subsection{Namespace discipline}
\label{sec:names}

Every name below is manufactured by quotation, not \texttt{new}-bound, so every
family is described by a predicate. $b$ is the board channel, $\mathit{pl}$ the
player channel, $g$ the game index.

\begin{center}
\begin{tabular}{@{}lll@{}}
\toprule
Family & Name & Payload \\
\midrule
board          & $\qt{[*\mathit{arena},g]}$                       & --- \\
occupancy      & $\qt{[*b,\texttt{"occ"},s]}$                     & $\Nil$ or $*p$ \\
attack         & $\qt{[*b,\texttt{"atk"},s,c,p]}$                 & attacker descriptor \\
pin            & $\qt{[*b,\texttt{"pin"},s]}$                     & pinning ray descriptor \\
check          & $\qt{[*b,\texttt{"chk"},c]}$                     & checking piece set \\
en passant     & $\qt{[*b,\texttt{"ep"},s]}$                      & capturable pawn channel \\
castling       & $\qt{[*b,\texttt{"cast"},c,\mathit{side}]}$      & right present or absent \\
piece          & $\qt{[*b,\texttt{"p"},\mathit{id}]}$             & --- \\
piece stack    & $\qt{[*b,\texttt{"m"},\mathit{id}]}$             & $\sigma$, the metabolic stack \\
advocacy       & $\qt{[*\mathit{pl},\texttt{"adv"},\qt{m}]}$      & advocating piece id \\
inhibition     & $\qt{[*\mathit{pl},\texttt{"inh"},\qt{m}]}$      & inhibiting piece id \\
franchise      & $\qt{[*\mathit{pl},\texttt{"V"},\mathit{id},r]}$ & the round's linear token \\
commit         & $\qt{[*\mathit{pl},\texttt{"commit"}]}$          & the chosen $\qt{m}$ \\
\bottomrule
\end{tabular}
\end{center}

\begin{requirement}[No family is receivable by the wrong society]
\label{req:sep}
No contract of the player receives on \texttt{occ} except through the ply join
of \texttt{Board}; no contract of the world receives on \texttt{adv},
\texttt{inh} or \texttt{V}. This is Proposition~\ref{prop:percept} discharged as
a code review checklist item, and it is the single most important invariant in
the system.
\end{requirement}

\subsection{The turn protocol}
\label{sec:protocol}

A turn proceeds as follows. Steps marked $\star$ are metered to the player;
steps unmarked are metered to the world.

\begin{enumerate}[leftmargin=2.2em]
\item \textbf{Settle.} \texttt{Settle} walks every ray from every occupied
square and deposits the \texttt{atk} family; then walks outward from each king
and deposits \texttt{pin} and \texttt{chk}. Terminating and confluent by
Remark~\ref{rem:settle}. On completion it signals the clock.
\item \textbf{Roll.} \texttt{Registry} publishes the set of live piece channels
for the moving colour.
\item \textbf{Franchise.} \texttt{Franchise} issues one linear $\mathsf{V}$
token per live piece into that piece's located stack. Round counter $r := 1$.
\item[$\star$] \textbf{Generate.} Each \texttt{Piece} enumerates its own
pseudo-legal moves from its own square by its own move contract. No piece
enumerates another's. The union over pieces is the candidate set; nobody
computes the union.
\item[$\star$] \textbf{Evaluate.} Each \texttt{Piece} evaluates its candidates
against the rung it currently holds, by querying \texttt{Check}. The query is an
assay: the formula sits in guard position, so evaluating is testing.
\item[$\star$] \textbf{Advocate.} Each \texttt{Piece} spends its $\mathsf{V}$
token on $\mathit{adv}(m)$ for its best candidate $m$, or declines.
\textsc{open}: whether a piece may advocate for a move it does not itself make
--- see \S\ref{sec:openspec}.
\item \textbf{Gate.} A \texttt{Gate} instance stands on each live candidate. The
first to accumulate $k$ advocacy sends fires and writes $\qt{m}$ to
\texttt{commit}. By Theorem~\ref{thm:2k} at most one fires when $k \geq
\lceil S/2 \rceil$.
\item \textbf{Round or resolve.} If no gate fired and $r < R$, then $r := r+1$
and control returns to step 3, with unconsumed advocacy left standing and
readable. If $r = R$, the clock applies
Definition~\ref{def:lastresort}.
\item \textbf{Ply.} \texttt{Board} performs the $N$-ary join of
Definition~\ref{def:ply}, guarded on the constraint family. A refused guard
leaves the message in the tuple space, which is the platform's own semantics for
a failed \texttt{where}, and the clock retries with the next-best candidate.
\item \textbf{Harvest.} If the join terminated an agent, its metabolic name is
routed per Definition~\ref{def:harvest}.
\item[$\star$] \textbf{Revise.} Each piece updates its rung from its own assay
outcomes by \texttt{Revise}. Refutation rate is the temperature, as in the
noughts world; it is an observable the assays already emit rather than a
schedule imposed from outside.
\item \textbf{Adjudicate.} \texttt{Arena} tests for mate, stalemate, and the
draw conditions, and either ends the game or passes the turn.
\end{enumerate}

\begin{remark}[Where the cost actually goes]
\label{rem:costgoes}
Steps 5 and 11 are the player's whole expenditure and step 1 is the world's.
That is the entire content of measurement M7. Note that step 4 is charged to the
player but is cheap and irreducible --- a piece must know its own moves --- and
step 7 is charged to the world because a join firing is a COMM, which under the
native integration is metered at the reducer on the deploy envelope.
\end{remark}

\subsection{The ladder}
\label{sec:ladder}

Each piece holds a rung. Rungs are cumulative and evaluated highest-priority
first, as in the noughts listing. $\mathsf{Emp}$, $\mathsf{Bears}$,
$\mathsf{Valuable}$, $\mathsf{Available}$, $\mathsf{Central}$ and
$\mathsf{Develops}$ are structural predicates over the deposited families;
$\KKmod{m}$ is the settled modality at the ply $m$.

\begin{center}
\begin{tabular}{@{}cl>{\raggedright\arraybackslash}p{6.1cm}c@{}}
\toprule
Rung & Character & Formula added & $\kappa$ \\
\midrule
0 & --- & $\top$ (any legal move) & $0$ \\
1 & non-modal & $\exists m.\ \mathsf{Capture}(m) \wedge \neg\,\mathsf{Defended}(\mathit{to}(m))$ & $2$ \\
2 & non-modal & $\vee\ \neg\,\mathsf{Hangs}(m)$: the destination is not attacked by a cheaper piece & $2$ \\
3 & spatial & $\vee\ \KKmod{m}\,\mathsf{Fork}(\mathit{me})$, Definition~\ref{def:fork} & $4$ \\
4 & alternating & $\vee\ \forall m'.\,\KKmod{m}\Bmod{m'}\,\neg\,\mathsf{Fork}(\mathit{them})$ & $8$ \\
5 & spatial, depth 0 & $\vee\ \mathsf{Central}(\mathit{to}(m)) \vee \mathsf{Develops}(m)$ & $1$ \\
6 & coalition & $\vee\ \KKmod{m}\,\mathsf{Bearing}(s) \geq v$, Definition~\ref{def:coalrung} & $16$ \\
$\chi$ & fixed point & $\nu X.\forall m.\Bmod{m}(\neg\,\mathsf{Lost} \wedge \exists m'.\KKmod{m'} X)$ & $32$ \\
\bottomrule
\end{tabular}
\end{center}

The prices are the noughts schedule $\kappa(d) = 2^d$ per candidate examined,
charged by the depth of the formula the rung installs, with rung~6 priced one
step above rung~4 because the coalition modality quantifies over opens and is a
width engine. Rung~5 is priced at~1 because it is name-level depth zero, and
that pricing is what made it the best marginal purchase in the noughts
measurement; whether it repeats here is M2.

\begin{remark}[$\chi$ is unreachable, and that is the point]
\label{rem:chi}
In noughts $\chi$ was exact at nine unfoldings because the board loses a square
every step. Chess has no such bound. $\chi$ is therefore listed for completeness
and as a limit object; no agent will hold it, and by Remark~\ref{rem:alive} that
is exactly the condition under which the ecology stays alive.
\end{remark}

\subsection{The cost schedule}
\label{sec:cost}

\begin{center}
\begin{tabular}{@{}lll@{}}
\toprule
Operation & Grade & Price \\
\midrule
read one deposited marker              & perception  & $1$ \\
evaluate a rung-$d$ formula, one candidate & perception & $\kappa(d) = 2^d$ \\
cached query, unchanged deposit        & perception  & $1$ \\
enumerate own pseudo-legal moves       & perception  & $1$ per square walked \\
spend an advocacy token                & interaction & $2$ \\
fire a commit gate                     & interaction & $2k$ \\
perform a ply join of arity $n$        & interaction & $4n$ \\
one ray step in \texttt{Settle}        & world        & $1$ \\
harvest                                & interaction & $0$ (routing, not work) \\
\bottomrule
\end{tabular}
\end{center}

Every price here is a \textsc{choice}, not a derivation, and the note should say
so in the same breath as reporting any number computed with it. The noughts work
varied its price scale over three settings and found the shape of its results
robust and the boundaries not; the same sensitivity analysis is required here and
is part of the acceptance criteria in \S\ref{sec:phases}.

\subsection{The checker interface}
\label{sec:checker}

The service assumed in \S\ref{sec:assumed} is used through one contract.

\begin{lstlisting}
// verdict in {"t","b","f"}; witness is Nil, a redex position, or a
// coalition (a list of piece ids); cost is charged to the caller's stack.
contract Check(@formula, @posHandle, @algebra, payer, ret) = {
  // ret!((verdict, value, witness, cost))
}
\end{lstlisting}

\begin{description}[leftmargin=1.6em,style=nextline]
\item[\texttt{formula}] a quoted formula of the generated logic extended with
$\Emod{N}{V}$, with namespaces given as name predicates.
\item[\texttt{posHandle}] a handle on the current deposit, not a copy of the
position. The checker reads the deposited families; it never drives the board.
This is Requirement~\ref{req:direction} enforced by the interface shape.
\item[\texttt{algebra}] one of \texttt{"B"}, \texttt{"Rnn"}, \texttt{"trop"},
\texttt{"vit"}. Witnesses are returned for the last two.
\item[\texttt{payer}] the caller's metabolic channel. A caller that cannot fund
the query receives verdict \texttt{"b"}, which is the budget bottom and is
deliberately indistinguishable from a stuck environment.
\end{description}

\begin{requirement}[The checker may not move]
\label{req:nomove}
The checker's implementation must receive only on the \texttt{atk},
\texttt{pin}, \texttt{chk}, \texttt{ep} and \texttt{cast} families and must
never receive on \texttt{occ}. Violating this makes the pride move the zebra by
looking at it, in the phrase the herd note used for the same hazard, and it is
the reason the noughts listing insisted on \texttt{peek} rather than a
consuming receipt in its own scan.
\end{requirement}

\subsection{Reproduction and the offline loop}
\label{sec:repro}

Between games the population evolves. This is the noughts world's
\texttt{Pick} / \texttt{Crossover} / \texttt{Beget} apparatus and it transfers
with one change of emphasis. There are two loci, the rung and the policy, and
the noughts measurement showed that an individual can reach the first by
deepening and cannot reach the second at all --- the opening changed only across
generations, by crossover, because it is a move of distance $2^{-1}$ in the
ultrametric and an individual's incremental revisions are confined to their
initial ball.

The chess analogue of ``the opening'' is the piece's \emph{advocacy style}: how
readily it advocates, what it treats as valuable, whether it defers to standing
advocacy. That is the locus crossover must reach, and it is where a model-backed
mutation operator earns its place (\S\ref{sec:policy}).

\subsection{Delivery phases}
\label{sec:phases}

\begin{description}[leftmargin=1.6em,style=nextline]
\item[P0 --- the world alone.] \texttt{Board}, \texttt{Settle},
\texttt{Registry}, \texttt{Clock}, \texttt{Arena}. Acceptance: legal move
generation validated against a standard perft suite to depth 5 from the initial
position and from a set of tactical positions; the settling shown terminating
and confluent by test; every ply a single join.
\item[P1 --- one agent, no society.] A single \texttt{Piece} contract driving
all pieces in turn, rung 0--2, no advocacy. Acceptance: plays legal complete
games against a fixed reference opponent; metering totals reconcile against
$\sigma + \kappa = \Theta$.
\item[P2 --- the society, arithmetic guards.] Sixteen agents, franchise, commit
gates, rounds, last resort. Rungs 0--2 and 5 only, since those are the non-modal
rungs and are expressible in the live guard fragment. Acceptance:
Theorem~\ref{thm:2k} confirmed empirically at $k$ below and above $S/2$;
last-resort firing rate reported.
\item[P3 --- the checker.] \texttt{Check} integrated; rungs 3, 4, 6 enabled.
Acceptance: M1 and M2 produced.
\item[P4 --- the economy.] Piece stacks, harvest, metabolic death, the offline
reproduction loop. Acceptance: M6 produced; conservation audited per ply.
\item[P5 --- the aperture variant.] $j$ of sixteen pieces backed by a model
session through a declared aperture, in a stratified shard. Acceptance: M8
produced; the static analysis of aperture channels demonstrated.
\end{description}

P0 through P2 are buildable today on
\texttt{feature/cost-accounted-rho}. P3 depends on the assumed checker. A
syntactically manipulable \texttt{Theory} of positions is \emph{not} on the
critical path and should not be put on it: the DDL and module syntax are absent
from the branch, which is the same dependency F1R3Skein is waiting on, and
nothing in P0--P5 requires them.

\subsection{Open specification items}
\label{sec:openspec}

\begin{description}[leftmargin=1.6em,style=nextline]
\item[S1] May a piece advocate for a move it does not itself make? Yes gives
genuine coalitions in the advocacy structure and lets a defender back an
attacker's plan; no keeps the franchise interpretation clean and makes advocacy
self-interested. Our inclination is yes, because it is what makes the advocacy
set a coalition rather than a poll, but it changes Theorem~\ref{thm:2k} not at
all and can be decided late.
\item[S2] Is $k$ absolute or a fraction of $S$? Proposition~\ref{prop:crossing}
says this determines whether the endgame is split-free or gate-dead.
\item[S3] Cross-inhibition in or out of the first build? Recommended out, per
\S\ref{sec:coord}.
\item[S4] Is the opponent a society too, or a conventional engine? Recommended
conventional for P1--P4: a fixed yardstick is needed, and the asymmetry is
itself the experiment. Society-versus-society belongs after M1.
\item[S5] Do pieces hold individual stacks from P2, or only from P4? Holding
them early makes the metering real but couples two sources of variance.
\item[S6] What is the reference opponent? It must be weak enough that the
society is not simply crushed and fixed enough that measurements are
comparable across phases.
\end{description}

\section{The measurements}
\label{sec:measure}

Eight measurements. Each states what would be reported and, where it matters,
what result would embarrass the framework.

\begin{description}[leftmargin=2.4em,style=nextline]
\item[M1 --- the yield ladder.] For each rung, win/draw/loss against a fixed
reference opponent at fixed strength, and cumulative metered cost, over a fixed
number of games. Reported as the noughts table was: win rate, cost, and marginal
return per metered unit. Unlike noughts these are sampled estimates with
confidence intervals, not exact counts, and the note reporting them must say so
in its own results section.
\item[M2 --- non-modal share.] The fraction of total yield delivered by rungs
1, 2 and 5 and the fraction of total cost they consume, against the noughts
figures of $92.3\%$ and $34.3\%$. This is the discriminating measurement of
Principle~\ref{prin:discriminate}. Either outcome is publishable and they say
opposite things.
\item[M3 --- the coalition rung.] Rung 6's marginal return, and separately
whether its witness is \emph{informative}: does the returned coalition coincide
with what a strong player would name? Scored by blind comparison on a set of
tactical positions. The herd note found that the modality bought irreversibility
rather than accuracy --- precision $1.0000$ at every sampled time, at low
coverage --- and if that repeats, the yield metric underprices instruments that
decline to answer and the second-currency principle acquires its second
instance.
\item[M4 --- split and deadlock against $k$.] Firing rate of the commit gate,
split rate, deadlock rate and last-resort rate, over $k$ from $2$ to $16$ at
$S = 16$. Theorem~\ref{thm:2k} predicts split rate identically zero for
$k \geq 9$, and the prediction is structural, so a single observed split above
that threshold is a bug in the franchise's linearity and should be treated as
such.
\item[M5 --- the endgame crossing.] Decision character as $S$ falls through
$2k$ within single games: rounds to commit, advocacy entropy, and last-resort
rate, bucketed by material remaining. This is the measurement
Remark~\ref{rem:predict} exists for and the one we would report first, because
it is the only prediction here that the source domain could not have made.
\item[M6 --- emergent values.] Lifetime harvest per piece type against the
classical $1:3:3:5:9$, per Conjecture~\ref{conj:values}. Report the raw ratios
and the regression, and report disagreement loudly if it occurs.
\item[M7 --- who is billed.] Total metered cost per game under the world-billed
and player-billed attack-map architectures of \S\ref{sec:billed}, split by
society. The framework predicts the world-billed architecture wins on total and
that the player's share falls by a factor near $S$. This is the cheapest
experiment in the set and it directly tests the read-versus-drive thesis.
\item[M8 --- the aperture's price.] With $j$ of sixteen pieces
model-backed, yield per metered unit as a function of $j$, against the
all-\flang\ baseline. Prices a model's judgement in the framework's currency.
\end{description}

\begin{remark}[The counter-move, if there is time]
\label{rem:counter}
The herd note's most striking result was the counter-move: an adversary that
introduces seams into the population does not induce \emph{error} in the
observer, it denies \emph{liveness} --- coverage collapsed from $0.7520$ to
$0.0115$ while precision stayed above $0.9986$. The chess analogue is an
opponent that plays to reduce the society's advocacy coherence rather than to
win material, and the prediction is that such an opponent should drive the
society to last resort rather than to blunders. If it holds, illegibility as
denial of liveness acquires an adversarial instance, and the legibility trade
--- that a population buys privacy with the same currency it buys agreement ---
becomes something one can demonstrate rather than argue.
\end{remark}

\section{What will go wrong}
\label{sec:risks}

\paragraph{Branching, and the honest expectation.} Chess has a branching factor
near thirty-five and rewards depth. A society that pays a rendezvous per
candidate cannot search, and the framework's own answer --- read the structure
the environment has deposited, do not drive it --- means this player will be a
perception player rather than a search player. It will lose to anything with a
search until its perception is very good. The evaluation criterion must
therefore be yield per metered unit, with playing strength as a secondary
sanity check, and this must be stated before the first game rather than
discovered in the results. A project of this kind fails publicly in exactly one
way, which is by having promised strength.

\paragraph{The society is small.} The swarm results are asymptotic in scout
population; at $S = 16$ the regimes are coarse and the advocacy distribution is
lumpy. This cuts both ways, since it is also what makes
Proposition~\ref{prop:crossing} reachable, but no fine-grained result about
quorum dynamics should be expected from sixteen agents.

\paragraph{The guard fragment.} Structural and modal guards are not executed by
the live reducer; they are specifications. P2 is buildable because rungs 1, 2
and 5 are expressible arithmetically over deposited markers, which is a
deliberate design of the ladder rather than a happy accident, but P3 is gated on
the assumed checker and the note should not pretend otherwise.

\paragraph{Latency.} Sixteen agents, up to $R$ rounds, a checker query per
candidate per agent. Even in the all-\flang\ configuration a turn is many
thousands of COMMs. A demonstration has to be watchable, so the round budget $R$
and the candidate cap per agent are the two dials that must be tuned early, and
they are both dials that \emph{change the science}, since capping candidates is
capping $\exists$. Report them with every number.

\paragraph{Reflection, a fifth time.} The subject-reduction hazard the programme
has now met four times is present here too: a received process can carry a
formula of higher altitude than the receiving clause declares, and this society
passes quoted moves and quoted formulae between agents constantly. The
stratified \texttt{where} discipline is the known repair and it should be
adopted from P2 rather than retrofitted, because retrofitting it means revisiting
every advocacy message.

\section{What is shown and what is not}
\label{sec:shown}

Shown, in the sense of proved here from imported premises:
Proposition~\ref{prop:percept} (perception must not act),
Proposition~\ref{prop:arity} with Corollary~\ref{cor:Ksurvives} (the join
taxonomy and the survival of $\Kmod{j}$), Proposition~\ref{prop:decay} and the
consequent choice of quorum algebra, and Theorem~\ref{thm:2k} with its
corollary that $k \geq \lceil S/2 \rceil$ buys structural no-split.
Theorem~\ref{thm:2k} is the only genuinely new theorem in the note and its proof
is three lines, which is the correct length: the work was done by choosing a
calculus in which linear tokens and $N$-ary joins are primitive.

Argued but not proved: that the world-billed architecture wins
(\S\ref{sec:billed}, awaiting M7); that legality belongs to the world
(Principle~\ref{prin:legality}, which is a design principle with a cost argument
behind it, not a theorem); that the model belongs outside the ply loop
(\S\ref{sec:policy}).

Conjectured: emergent piece values (Conjecture~\ref{conj:values}).

Not addressed at all: opening theory and endgame tablebases, both of which are
external knowledge and both of which the society has no principled way to
consume; the fifty-move and repetition rules, which are history-dependent and
therefore want the history monad rather than the deposit; and time control,
which would make the budget a clock and is probably the most interesting thing
omitted.

\section{Open problems}
\label{sec:open}

\begin{description}[leftmargin=2.4em,style=nextline]
\item[Q1] Does the noughts yield result survive? This is M2 and it is the
project's reason for existing.
\item[Q2] Is there a principled $k$? The design fixes $k \geq \lceil S/2
\rceil$ for safety, but nothing says which value in that range is best, and the
swarm's answer --- that $k$ is not a number but a rate threshold, scaling with
encounter geometry --- suggests the chess analogue is a rate as well.
\item[Q3] Does agent-per-plan beat agent-per-piece? Remark~\ref{rem:decomp}.
The obstacle is that a plan cannot be captured, so predation would have to be
redefined.
\item[Q4] What is the right treatment of the fifty-move and repetition rules?
They are the only genuinely historical facts in chess and the programme has a
history monad that nothing in this design uses.
\item[Q5] Does the coalition modality's irreversibility-not-accuracy signature
repeat? M3.
\item[Q6] Can the society learn its own \texttt{Open}? The herd note left this
open in general; in chess the analogue is whether the pieces can discover which
subsets of themselves are worth treating as coalitions, rather than having the
namespace predicate handed to them.
\item[Q7] Does subject reduction hold for the stratified discipline in the
presence of quoted moves passed as advocacy payloads? The fifth appearance of
the same suspect.
\item[Q8] Is there a composition law for societies? Two coalitions that each
bear on a square are not obviously composable, and the composing-learners note's
disjointness-versus-overlap dial is the obvious place to look.
\item[Q9] What does time control do? Making the budget a clock rather than a
stack changes the metabolism from conservation to flow, and it is not obvious
the conservation results survive.
\item[Q10] Does the society's play have a signature a human can recognise? Not
a formal question, but the one a reader will ask, and the witness coalitions of
Corollary~\ref{cor:witness} are the only handle on it.
\end{description}

\section{Conclusion}
\label{sec:conclusion}

The proposal was to make each piece an agent. The analysis above says that this
is a well-formed thing to do in cost-accounted \rhoc, that three of its
apparently difficult parts are already solved by constructions in hand, and that
one of them is easier here than in the setting it came from.

The board encoding wants three name families rather than one, because perception
must not be able to act and because the attack map is the shared expensive
object whose billing is the project's cheapest experiment. A ply is an $N$-ary
join and the join arity is the move taxonomy, which is what keeps the generated
modality generated. Coordination is quorum sensing rather than consensus, and the
argument for that is a derivation from budget rather than a design preference.
Legality is the one unavoidable universal and it belongs to the world, which has
the pleasant consequence that pin-awareness becomes perception.

The part that is easier here than in the swarm is the two-quorum theorem. In the
bees it is a structural impossibility result confirmed by simulation. In
\rhoc\ it is arithmetic: a quorum is a join arity, one linear franchise token
per piece per round is the polity's dereliction, and two $k$-ary joins over a
pool of $S$ linear tokens need $2k \leq S$. A designer who sets $k \geq
\lceil S/2\rceil$ has bought a safety property outright, for any policies
whatsoever, and has thereby confined cross-inhibition to the single job of
breaking deadlock, where a cheaper mechanism --- reading the standing advocacy
between rounds --- is available.

What the project is for is narrower than what it will look like. It will look
like a chess engine and it will not be a good one. It is an instrument for one
question: whether the finding that non-modal rungs deliver most of the yield for
a third of the cost is a fact about games or a fact about a game small enough to
solve. Noughts-and-crosses cannot answer that, because everything about it is
exact and nothing about it is hard. Chess can, and the same board that answers
it also happens to be a world where the ecology cannot die of its own
perfection.

\appendix

\section{The listing}
\label{app:rho}

Conventions follow the game note. A parameter written \texttt{c} is a NAME; a
parameter written \texttt{@d} is DATA; a channel is passed as \texttt{*c}, since
\texttt{@(*c) = c}. Guards marked \texttt{//! ideal} use structural or modal
connectives that the current reducer does not evaluate; they are specifications
of a check rather than checks. Everything else runs. The listing is the spine
only: move generation per piece type, the full settling walk and adjudication are
elided, since they are long and unsurprising, and the file
\texttt{chess.rho} accompanying this note carries them.

\lstinputlisting{chess.rho}

\section{The formula catalogue}
\label{app:formulae}

Structural predicates, all over deposited families, all at perception prices:

\begin{align*}
  \mathsf{Occ}(s,p) &:= \out(\mathit{occ}(s))\,\qt{(*p)} \\
  \mathsf{Emp}(s) &:= \out(\mathit{occ}(s))\,\qt{\Nil} \\
  \mathsf{Bears}(p,s) &:= \exists c.\ \out(\mathit{atk}(s,c,p))\,\top \\
  \mathsf{Attacked}(s,c) &:= \exists p.\ \mathsf{Bears}(p,s) \wedge \mathsf{Colour}(p,c) \\
  \mathsf{Defended}(s) &:= \mathsf{Attacked}(s,\mathit{me}) \\
  \mathsf{Pinned}(p) &:= \exists s.\ \mathsf{Occ}(s,p) \wedge \out(\mathit{pin}(s))\,\top \\
  \mathsf{Available}(p) &:= \neg\,\mathsf{Pinned}(p) \wedge \neg\,\mathsf{Committed}(p) \\
  \mathsf{InCheck}(c) &:= \out(\mathit{chk}(c))\,\top
\end{align*}

Derived formulae:

\begin{align*}
  \mathsf{Threat}(p) &:= \exists s.\ \mathsf{Bears}(p,s) \wedge \mathsf{Valuable}(s) \\
  \mathsf{Fork}(p) &:= \mathsf{Threat}(p) \gsep{0}{N_{\mathit{atk}}} \mathsf{Threat}(p) \\
  \mathsf{Hangs}(m) &:= \mathsf{Attacked}(\mathit{to}(m),\mathit{them}) \wedge
                        \neg\,\mathsf{Defended}(\mathit{to}(m)) \\
  \mathsf{Bearing}(s) &:= \Emod{N_{\mathit{pl}}}{V}\ \lambda p.\
                        \mathsf{Bears}(p,s) \otimes \mathsf{Available}(p) \\
  \mathsf{Coherent} &:= \nu X.\ \Emod{N_{\mathit{pl}}}{V}\big(\lceil \beta \rceil
                        \otimes \langle K \rangle X\big)
\end{align*}

$\mathsf{Coherent}$ is the society's analogue of herd coherence: the coalition
keeps agreeing as the population steps. It is not on the ladder because it is a
$\nu$ over the run tree and therefore a depth engine, priced accordingly; it is
listed because it is the formula an \emph{observer} of the society would use,
and the society is going to be observed.

\end{document}
